% Source: https://artofproblemsolving.com/community/c3h1335921_amc_8_2007_problem_17
% Stable identity: aops-topic-1335921-q01
The original mixture contains $30\%\cdot30=9$ liters of yellow tint. After adding $5$ liters there are $14$ liters of yellow tint in $35$ liters of paint, so the percentage is
\[
\dfrac{14}{35}=\dfrac{2}{5}=40\%,\qquad\mathrm{(C)}\ \boxed{40}.
\]
(Alternative: the non-yellow part, $70\%\cdot30=21$ liters, is unchanged and is now $\dfrac{21}{35}=60\%$ of the mixture, so yellow is $40\%$.)
